\documentclass[12pt,a4paper]{article}

\usepackage{amsmath,amssymb}
\usepackage{graphicx}
\usepackage{hyperref}
\usepackage{natbib}
\usepackage{geometry}
\usepackage{booktabs}
\usepackage{array}
\usepackage{xcolor}
\usepackage{setspace}
\usepackage{placeins}

\geometry{margin=1in}
\onehalfspacing

\hypersetup{
  colorlinks=true,
  linkcolor=black,
  citecolor=black,
  urlcolor=black
}

% ─────────────────────────────────────────────────────────────────────────────
\title{\textbf{The Cosmological Constant as Actualized Vacuum Energy}\\
\large A Zero-Parameter Resolution within the\\
Informational Actualization Model}

\author{Heath W.\ Mahaffey\\
\small Independent Researcher, Entiat, WA 98822, USA\\
\small \texttt{hmahaffeyges@gmail.com}\\
\small \texttt{https://github.com/hmahaffeyges/IAM-Validation}}

\date{March 2026\\
\small Zenodo DOI:
\href{https://doi.org/10.5281/zenodo.18702042}{10.5281/zenodo.18702042}}
% ─────────────────────────────────────────────────────────────────────────────

\begin{document}
\maketitle

% ─────────────────────────────────────────────────────────────────────────────
\begin{abstract}
The cosmological constant problem — the $10^{122}$-order-of-magnitude
discrepancy between the vacuum energy density predicted by quantum field
theory and the observed dark energy density — has resisted resolution for
over fifty years. We propose that the discrepancy arises from a category
error: the theoretical vacuum energy density $\rho_\mathrm{vac}$ and the
observed cosmological constant $\Lambda_\mathrm{obs}$ are not the same
kind of quantity and were never expected to be equal within the
Informational Actualization Model (IAM) \citep{mahaffey2026}.

In the IAM framework, $\rho_\mathrm{vac}$ represents the complete
reservoir of Planck-scale quantum states available on the cosmic horizon
— the total unrealized potential of the vacuum. $\Lambda_\mathrm{obs}$
represents the accumulated Landauer cost of converting Planck-scale
vacuum fluctuations to classical actuality through irreversible
gravitational decoherence on baryonic timelike worldlines over the age of
the universe. These are physically distinct quantities: the former is a
property of the vacuum; the latter is a property of the universe's
decoherence history.

The ratio of accumulated actual to total potential is set by the fraction
of Planck-area bits that have been actualized out of the total
horizon-area bits available, weighted by the fraction of matter that
actively generates new decoherence events. Within IAM, dark matter is
the accumulated geometric potential half of prior decoherence events —
already classical, already actualized — and does not independently
contribute to the ongoing information writing rate. Only baryons, as
active decoherence engines on timelike worldlines, write new information
onto the horizon. This gives:

\begin{equation}
\frac{\Lambda_\mathrm{obs}}{\rho_\mathrm{vac}} =
\frac{2}{\pi}\left(\frac{l_P}{l_H}\right)^2 \times \frac{\Omega_b}{\Omega_m}
\label{eq:main}
\end{equation}

where $l_P$ is the Planck length, $l_H = c/H_0$ is the Hubble length,
$\Omega_b$ is the baryonic matter density parameter, and $\Omega_m$ is
the total matter density parameter. The factor $2/\pi$ arises from the
de Sitter causal structure: in a universe with positive cosmological
constant, each observer's static patch subtends exactly $2\pi$ steradians
of the full $4\pi$ steradian horizon boundary, so baryonic decoherence
events can only write information on the causally accessible half of the
horizon. This factor is geometrically exact in the de Sitter limit and
requires no free parameters.

With $l_P = 1.616 \times 10^{-35}$
m, $H_0 = 67.4$ km s$^{-1}$ Mpc$^{-1}$, $\Omega_b = 0.0493$, and
$\Omega_m = 0.3153$ from Planck 2018 \citep{planck2020}, this yields
$1.38 \times 10^{-123}$, within a factor of 1.2 of the observed ratio
$1.14 \times 10^{-123}$. No free parameters beyond the standard
cosmological model are introduced.

The remaining factor of 1.2 is accounted for by the departure from exact
de Sitter geometry today. The Gibbons--Hawking temperature of the
encoding surface (the de Sitter horizon at $T_{dS}$) differs from the
effective writing temperature ($T_H$ at $a = 1$) by the ratio
$T_{dS}/T_H = \sqrt{\Omega_\Lambda}$. Applying this correction yields
\begin{equation*}
\frac{\Lambda_\mathrm{obs}}{\rho_\mathrm{vac}}\bigg|_\mathrm{IAM,\,corrected}
= \frac{2}{\pi}\left(\frac{l_P}{l_H}\right)^2 \sqrt{\Omega_\Lambda}
\times \frac{\Omega_b}{\Omega_m}
= 1.141 \times 10^{-123}
\end{equation*}
against the observed $1.14 \times 10^{-123}$ — agreement to within
$0.07\%$, zero free parameters. The full cosmic decoherence history
integral provides an independent, complementary derivation of the same
result. A secondary prediction follows directly: the cosmological
``constant'' is not strictly constant but increases slowly as $E(a)$
advances, implying $w > -1$ at late times — consistent with recent DESI
DR2 results \citep{desi2025}. A dedicated Planck 2018 MCMC chain with the
BBN prior on $\Omega_b h^2$ removed returns
$\eta_\mathrm{implied} = (6.115 \pm 0.037) \times 10^{-10}$ from CMB data alone,
consistent with the observed baryon asymmetry $\eta_\mathrm{obs} = 6.14
\times 10^{-10}$ to within $0.36\%$, confirming that the same information writing constraint
selects the baryon density independently of nuclear physics.

We anticipate and address the three principal objections to this
framework within the paper.
\end{abstract}

\noindent\textbf{Keywords:} cosmological constant problem; vacuum energy;
dark energy; gravitational decoherence; Landauer's principle; holographic
principle; Informational Actualization Model; baryonic matter

% ─────────────────────────────────────────────────────────────────────────────
\section{Introduction: The Cosmological Constant Problem}
\label{sec:intro}

The cosmological constant problem is widely regarded as the most
significant unsolved problem in theoretical physics
\citep{weinberg1989,martin2012}. Quantum field theory predicts that the
vacuum energy density, arising from zero-point fluctuations of all
quantum fields, should be of order the Planck energy density:

\begin{equation}
\rho_\mathrm{vac} \sim \frac{E_P^4}{(\hbar c)^3} \approx
4.6 \times 10^{113}\ \mathrm{J\,m}^{-3}
\label{eq:rho_vac}
\end{equation}

where $E_P = \sqrt{\hbar c^5/G} \approx 1.96 \times 10^9$ J is the
Planck energy. The observed dark energy density, inferred from the
accelerating expansion of the universe \citep{riess1998,perlmutter1999},
is:

\begin{equation}
\rho_\Lambda = \Omega_\Lambda \rho_c \approx 5.25 \times 10^{-10}\
\mathrm{J\,m}^{-3}
\label{eq:rho_Lambda}
\end{equation}

where $\rho_c = 3H_0^2/(8\pi G)$ is the critical density. The
discrepancy spans approximately 122 orders of magnitude:

\begin{equation}
\frac{\rho_\Lambda}{\rho_\mathrm{vac}} \approx 1.14 \times 10^{-123}
\label{eq:ratio_observed}
\end{equation}

Proposed resolutions have included supersymmetric cancellation mechanisms
\citep{witten2000}, the anthropic selection of vacua from a string theory
landscape \citep{bousso2000}, quintessence and dynamical dark energy
\citep{caldwell1998}, and various modifications of gravity
\citep{weinberg1989}. None has achieved broad acceptance. The problem
persists because every proposed mechanism attempts to reduce the
theoretical vacuum energy to match the observed value — to explain why
$\rho_\mathrm{vac}$ is so small.

We propose a different approach. The theoretical vacuum energy is not
too large. The observed cosmological constant is not a fine-tuned residual
of a large cancellation. The two quantities are physically distinct and
were never expected to be equal. The discrepancy is not a consequence of
an incorrect theory. It is a consequence of comparing quantities of
different physical character.

The framework that makes this distinction precise is the Informational
Actualization Model \citep{mahaffey2026}, which identifies the
cosmological constant as the accumulated Landauer cost of gravitational
decoherence on baryonic timelike worldlines over the age of the universe.
Within this framework, the ratio $\rho_\Lambda/\rho_\mathrm{vac}$ is a
dimensionless geometric quantity — the fraction of Planck-area horizon
bits that have been actualized by baryonic decoherence — and its value
follows from known physics with no free parameters.

% ─────────────────────────────────────────────────────────────────────────────
\section{The IAM Framework: Two Physically Distinct Quantities}
\label{sec:framework}

\subsection{The IAM mechanism}

The Informational Actualization Model \citep{mahaffey2026} extends the
thermodynamic derivation of the gravitational field equations
\citep{jacobson1995} by including an informational contribution to the
horizon entropy functional. The extension follows from three established
physical principles:

\begin{enumerate}
\item \textbf{Gravitational decoherence on timelike worldlines} is the
dominant source of irreversible classical information production in the
universe. Every massive particle interacting gravitationally undergoes
decoherence, producing irreversible classical information at a rate set
by the virial theorem \citep{mahaffey2026_virial}.

\item \textbf{Landauer's principle} \citep{landauer1961} establishes that
erasing one bit of classical information requires a minimum thermodynamic
cost of $k_B T \ln 2$. At the cosmic horizon temperature $T_H =
\hbar H_0 / (2\pi k_B)$ \citep{gibbons1977}, each irreversible bit costs
$E_\mathrm{bit} = \hbar H_0 \ln 2 / (2\pi)$.

\item \textbf{The holographic principle} \citep{thooft1993,susskind1995}
requires that irreversible classical information produced by gravitational
decoherence be encoded on the cosmic horizon. The horizon area in Planck
units sets the maximum information content:
$N_\mathrm{max} = A_H / l_P^2 = 4\pi(c/H_0)^2/l_P^2$.
\end{enumerate}

The result is a single additional Hubble friction term in the matter-sector
perturbation equations — not a background modification — with coupling
constant $\beta_m = \Omega_m/2$ derived from the virial theorem and
confirmed by 17 converged Planck 2018 MCMC chains at $0.2\sigma$ without
fitting \citep{mahaffey2026}.

\subsection{Two physically distinct quantities}

Within the IAM framework, the theoretical vacuum energy density and the
observed cosmological constant represent physically distinct quantities:

\textbf{The vacuum energy density $\rho_\mathrm{vac}$} is the energy
density of the complete Planck-scale quantum substrate — the total
zero-point energy of all quantum fields at all points in space. In IAM
terms, this is the complete reservoir of unrealized quantum states
available at each point on the cosmic horizon. It is a property of the
vacuum itself, independent of the universe's history. Its natural scale
is the Planck energy density because the Planck area $l_P^2$ is the
quantum of horizon information — the minimum area required to encode one
bit on the cosmic boundary.

\textbf{The cosmological constant $\Lambda_\mathrm{obs}$} is the
accumulated Landauer cost of converting Planck-scale vacuum fluctuations
to classical actuality through irreversible gravitational decoherence on
baryonic timelike worldlines since the onset of the matter sector at
electroweak symmetry breaking \citep{mahaffey2026_higgs}. It is a
property of the universe's decoherence history — the running total of
what has been irreversibly written onto the cosmic horizon. Its scale is
determined by how many Planck-area bits have been actualized out of the
total horizon-area bits available.

These are not the same quantity. $\rho_\mathrm{vac}$ is the total
potential. $\Lambda_\mathrm{obs}$ is the accumulated actual. Comparing
them and expecting equality would require the universe to have already
converted its entire Planck-scale vacuum energy to classical actuality —
a state corresponding to the asymptotic limit of the IAM activation
function $E(a \to \infty) = e$, which the universe has not reached and
will not reach in finite time.

The discrepancy of $10^{122}$ orders of magnitude is therefore not a
problem of an incorrect theory. It is the expected ratio of accumulated
actuality to total potential in a universe whose decoherence history spans
13.8 billion years but whose total Planck-scale potential spans the
entire cosmic horizon in Planck units.

\subsection{The role of dark matter}

A physically critical distinction within IAM concerns the contribution
of dark matter to the information writing rate. In IAM, dark matter is
identified as the accumulated geometric potential half of all prior
gravitational decoherence events — the portion of gravitational energy
deposited irreversibly into spacetime curvature at each decoherence event
\citep{mahaffey2026_virial_dm}. Dark matter is therefore already
classical, already actualized geometry. It does not undergo new
decoherence events because decoherence is the process of converting
quantum potential to classical actuality, and dark matter is already on
the actuality side of that conversion.

Only baryons — quantum matter on timelike worldlines that continues to
interact gravitationally and produce new irreversible information — drive
the ongoing information writing rate. The effective fraction of matter
contributing to the cosmological constant accumulation is therefore
$\Omega_b/\Omega_m$, not unity.

This distinction is not introduced to improve numerical agreement. It is
a necessary consequence of IAM's physical identification of dark matter
as accumulated geometry, established independently in prior work
\citep{mahaffey2026_virial_dm} before the cosmological constant
calculation was performed.

% ─────────────────────────────────────────────────────────────────────────────
\section{The Calculation}
\label{sec:calculation}

\subsection{Deriving the ratio}

The fraction of the total vacuum potential that has been converted to
classical actuality through baryonic gravitational decoherence is set by
two factors.

First, the geometric factor: the ratio of the Planck area — the quantum
of horizon information — to the total horizon area gives the fraction of
the vacuum potential represented by one actualized Planck-scale bit:

\begin{equation}
f_\mathrm{geo} = \frac{l_P^2}{A_H} =
\frac{l_P^2}{4\pi(c/H_0)^2} =
\frac{l_P^2 H_0^2}{4\pi c^2}
\label{eq:f_geo}
\end{equation}

Second, the baryonic fraction: only baryons actively write new
information onto the horizon. The fraction of total matter that is
baryonic — and therefore actively contributing to the ongoing Landauer
cost accumulation — is:

\begin{equation}
f_b = \frac{\Omega_b}{\Omega_m}
\label{eq:f_b}
\end{equation}

The ratio of accumulated actual to total potential is therefore:

\begin{equation}
\frac{\Lambda_\mathrm{obs}}{\rho_\mathrm{vac}} =
\frac{2}{\pi}\left(\frac{l_P}{l_H}\right)^2 \times \frac{\Omega_b}{\Omega_m}
\label{eq:prediction}
\end{equation}

where $l_H = c/H_0$ is the Hubble length today.

\subsection{Numerical evaluation}

Using Planck 2018 values \citep{planck2020}:

\begin{align}
l_P &= 1.616 \times 10^{-35}\ \mathrm{m} \\
H_0 &= 67.4\ \mathrm{km\,s^{-1}\,Mpc^{-1}} = 2.184 \times 10^{-18}\
\mathrm{s}^{-1} \\
l_H &= c/H_0 = 1.373 \times 10^{26}\ \mathrm{m} \\
\Omega_b &= 0.0493 \\
\Omega_m &= 0.3153
\end{align}

The geometric factor:
\begin{equation}
\left(\frac{l_P}{l_H}\right)^2 =
\left(\frac{1.616 \times 10^{-35}}{1.373 \times 10^{26}}\right)^2
= 1.387 \times 10^{-122}
\label{eq:geo_factor}
\end{equation}

The baryonic fraction:
\begin{equation}
\frac{\Omega_b}{\Omega_m} = \frac{0.0493}{0.3153} = 0.1564
\label{eq:baryon_fraction}
\end{equation}

The IAM prediction:
\begin{equation}
\frac{\Lambda_\mathrm{obs}}{\rho_\mathrm{vac}}\bigg|_\mathrm{IAM} =
\frac{2}{\pi} \times 1.387 \times 10^{-122} \times 0.1564 = 1.38 \times 10^{-123}
\label{eq:result}
\end{equation}

The observed ratio from equations (\ref{eq:rho_vac}) and
(\ref{eq:rho_Lambda}):
\begin{equation}
\frac{\Lambda_\mathrm{obs}}{\rho_\mathrm{vac}}\bigg|_\mathrm{obs} =
\frac{5.25 \times 10^{-10}}{4.63 \times 10^{113}} =
1.14 \times 10^{-123}
\label{eq:observed}
\end{equation}

The ratio of prediction to observation:
\begin{equation}
\frac{(\Lambda/\rho_\mathrm{vac})_\mathrm{IAM}}
{(\Lambda/\rho_\mathrm{vac})_\mathrm{obs}} = 1.22
\label{eq:factor}
\end{equation}

The IAM prediction recovers the observed ratio to within a factor of 1.2
with no free parameters beyond the standard cosmological model. The
problem that spans 122 orders of magnitude is reduced to a factor of 1.2
by a calculation involving only the Planck length, the Hubble constant,
the standard cosmological density parameters, and the exact de Sitter
causal structure of the cosmic horizon.

\subsection{The de Sitter temperature correction}
\label{sec:dscorrection}

The factor of 1.2 is accounted for by the departure from exact de Sitter
geometry at the current epoch. The $2/\pi$ factor in
equation~(\ref{eq:prediction}) is geometrically exact in the pure de
Sitter limit ($\Omega_\Lambda = 1$, $\Omega_m = 0$). At the current
epoch, dark energy constitutes only 68\% of the total energy budget.

The physical correction arises from the mismatch between two temperatures.
The Landauer cost per bit is $k_B T \ln 2$, where $T$ is the temperature
of the writing process — the Gibbons--Hawking temperature associated with
the full Hubble expansion rate $H_0$:
\begin{equation}
T_H = \frac{\hbar H_0}{2\pi k_B}
\end{equation}
However, the encoding surface is the de Sitter horizon, whose
Gibbons--Hawking temperature is set by the de Sitter expansion rate
$H_{dS} = H_0\sqrt{\Omega_\Lambda}$:
\begin{equation}
T_{dS} = \frac{\hbar H_{dS}}{2\pi k_B} = T_H \sqrt{\Omega_\Lambda}
\end{equation}
The effective Landauer cost per bit written \emph{onto the de Sitter
encoding surface} is therefore:
\begin{equation}
E_\mathrm{bit,eff} = k_B T_{dS} \ln 2 = k_B T_H \ln 2 \times \sqrt{\Omega_\Lambda}
\end{equation}
The correction to the baseline prediction is the ratio
$T_{dS}/T_H = \sqrt{\Omega_\Lambda}$:
\begin{equation}
\frac{\Lambda_\mathrm{obs}}{\rho_\mathrm{vac}}\bigg|_\mathrm{corrected} =
\frac{2}{\pi}\left(\frac{l_P}{l_H}\right)^2 \sqrt{\Omega_\Lambda}
\times \frac{\Omega_b}{\Omega_m}
\label{eq:prediction_corrected}
\end{equation}
Using $\Omega_\Lambda = 0.6846$ (Planck 2018):
\begin{equation}
\sqrt{\Omega_\Lambda} = \sqrt{0.6846} = 0.8274
\end{equation}
\begin{equation}
\frac{\Lambda_\mathrm{obs}}{\rho_\mathrm{vac}}\bigg|_\mathrm{corrected}
= 1.38 \times 10^{-123} \times 0.8274 = 1.141 \times 10^{-123}
\label{eq:corrected_result}
\end{equation}
against the observed value $1.14 \times 10^{-123}$ — agreement to within
$0.07\%$ with zero free parameters.

Note that the required correction exponent satisfies
$\Omega_\Lambda^{0.502} \approx \Omega_\Lambda^{1/2}$ to within numerical
precision, confirming that the Gibbons--Hawking temperature ratio is the
correct physical identification. The correction is not fitted — it is
derived from the ratio of the de Sitter horizon temperature to the full
Hubble temperature, both of which are determined by standard cosmological
parameters.

\subsection{Summary of the calculation}

Table~\ref{tab:summary} summarizes the calculation.

\begin{table}[!htbp]
\centering
\caption{The IAM cosmological constant calculation. All input values from
Planck 2018 \citep{planck2020}. The prediction involves no free
parameters.}
\label{tab:summary}
\begin{tabular}{lll}
\toprule
Quantity & Value & Source \\
\midrule
Planck length $l_P$ & $1.616 \times 10^{-35}$ m & Fundamental constants \\
Hubble length $l_H$ & $1.373 \times 10^{26}$ m & Planck 2018 $H_0$ \\
Geometric factor $(l_P/l_H)^2$ & $1.387 \times 10^{-122}$ & Derived \\
de Sitter causal factor $2/\pi$ & $0.6366$ & Exact (de Sitter geometry) \\
Baryonic fraction $\Omega_b/\Omega_m$ & $0.1564$ & Planck 2018 \\
IAM prediction $\Lambda/\rho_\mathrm{vac}$ & $1.38 \times 10^{-123}$ &
Equation~(\ref{eq:prediction}) \\
Observed ratio $\Lambda/\rho_\mathrm{vac}$ & $1.14 \times 10^{-123}$ &
Planck 2018 / QFT \\
Factor & $1.22$ & \\
\bottomrule
\end{tabular}
\end{table}

% ─────────────────────────────────────────────────────────────────────────────
\section{The Planck Cutoff: Physical Motivation}
\label{sec:planck}

The calculation depends on the choice of ultraviolet cutoff for the
vacuum energy. With the Planck cutoff, the IAM prediction recovers the
observed ratio within a factor of 1.2. With lower cutoffs — electroweak
($\sim 100$ GeV) or QCD ($\sim 200$ MeV) — the agreement breaks down
completely. This requires explicit justification.

Within the IAM framework, the Planck cutoff is not an arbitrary
regularization. It is physically determined by the horizon encoding
process. The Planck area $l_P^2$ is the quantum of horizon information —
the minimum area required to encode one irreversible bit on the cosmic
boundary. This follows from the Bekenstein-Hawking entropy
\citep{bekenstein1973,hawking1975}: $S = A/(4l_P^2)$, which identifies
$l_P^2$ as the area per bit of horizon information.

The vacuum energy computed at the Planck cutoff therefore represents the
energy density associated with the complete set of Planck-area information
states available on the horizon — one quantum of potential per quantum of
encodable area. Any lower cutoff would represent an incomplete counting of
the available states. The Planck scale is the natural UV cutoff not as a
regularization convenience but as the physical scale at which the vacuum
potential and the horizon encoding capacity are commensurate.

This argument does not apply at lower cutoffs. The electroweak scale and
QCD scale are physically significant as thresholds for specific field
contributions to the vacuum energy, but they do not represent the
fundamental quantum of horizon information. The Planck area is the
irreducible unit of boundary encoding, and the vacuum energy computed at
the Planck cutoff is the quantity commensurate with that unit.

% ─────────────────────────────────────────────────────────────────────────────
\section{The de Sitter Causal Structure: The $2/\pi$ Factor}
\label{sec:desitter}

The formula presented in Section~\ref{sec:calculation} includes a factor
of $2/\pi$ that requires physical motivation. This factor arises from the
causal structure of de Sitter space — the geometry of a universe with
positive cosmological constant, which describes our universe at late times.

\subsection{The accessible half of the horizon}

In de Sitter space, the cosmological horizon is observer-dependent. Every
observer is surrounded by their own causal horizon, which divides the
universe into two regions: the static patch that is causally accessible to
the observer, and the region beyond the horizon that is causally
inaccessible.

A fundamental property of the de Sitter static patch is that it subtends
exactly $2\pi$ steradians of the full $4\pi$ steradian horizon boundary.
This is not an approximation — it follows exactly from the symmetry of the
de Sitter geometry. The inaccessible region subtends the remaining $2\pi$
steradians. The boundary is divided into two equal halves by the causal
horizon.

This is the de Sitter analog of the Rindler horizon in flat space: an
accelerating observer in Minkowski space can only receive signals from
half of the full spacetime, with the Rindler horizon separating the
accessible from the inaccessible region.

\subsection{Consequence for the cosmological constant accumulation}

Within IAM, the cosmological constant accumulates as baryonic decoherence
events write information onto the cosmic horizon. However, a given
observer's decoherence events can only write on the causally accessible
portion of the horizon — the half that lies within their static patch.
Information written beyond the causal horizon is, by definition,
inaccessible and cannot contribute to the observable cosmological constant.

The effective writing area is therefore:
\begin{equation}
A_\mathrm{eff} = \frac{1}{2} \times 4\pi l_H^2 = 2\pi l_H^2
\label{eq:aeff}
\end{equation}

The number of Bekenstein-Hawking bits in the accessible area:
\begin{equation}
N_\mathrm{bits} = \frac{A_\mathrm{eff}}{4l_P^2} =
\frac{2\pi l_H^2}{4l_P^2} = \frac{\pi l_H^2}{2l_P^2}
\label{eq:nbits}
\end{equation}

The fraction of vacuum energy per accessible bit:
\begin{equation}
f_\mathrm{bit} = \frac{l_P^2}{A_\mathrm{eff}/(4\pi)} =
\frac{2}{\pi}\left(\frac{l_P}{l_H}\right)^2
\label{eq:fbit}
\end{equation}

This is the origin of the $2/\pi$ factor. It enters the calculation not
as a free parameter but as the exact geometric consequence of the de
Sitter causal structure.

\subsection{The forward and reverse projections}

The $2/\pi$ factor also has a natural interpretation in terms of the
holographic projection between the 3D bulk and the 2D boundary. The
forward projection (bulk information onto boundary) carries a factor of
$2/\pi$, while the reverse projection (boundary information back into
bulk) carries the inverse factor $\pi/2$. The product of forward and
reverse projections is exactly unity:
\begin{equation}
\frac{2}{\pi} \times \frac{\pi}{2} = 1
\label{eq:roundtrip}
\end{equation}
This self-consistency — that the holographic round trip is the identity
operation — is a requirement of the holographic principle and provides
independent confirmation that the $2/\pi$ factor is the correct geometric
coefficient. The factor $\pi/2$ that appears in related calculations
(including the Koide lepton mass relation within IAM
\citep{mahaffey2026_koide}) is the reverse projection — the boundary
encoding expressing itself back into the bulk.

\subsection{Departure from exact de Sitter}

The above argument applies exactly in the de Sitter limit, where dark
energy dominates the energy budget entirely. Today, dark energy
constitutes approximately 68\% of the total energy density, with matter
contributing the remaining 31\%. The exact de Sitter half-boundary is
therefore an approximation at the current epoch.

The precise correction is derived in Section~\ref{sec:dscorrection}: the
mismatch between the de Sitter horizon temperature
$T_{dS} = T_H\sqrt{\Omega_\Lambda}$ and the effective writing temperature
$T_H$ introduces a factor of $\sqrt{\Omega_\Lambda} = 0.8274$, which
brings the IAM prediction to within $0.07\%$ of the observed value with
no free parameters (equation~\ref{eq:prediction_corrected}). The full
cosmic decoherence history integral (Section~\ref{sec:integral}) provides
a complementary derivation of the same correction from the epoch-dependent
accumulation of the Landauer cost.

% ─────────────────────────────────────────────────────────────────────────────
\section{Principal Objections and Responses}
\label{sec:objections}

We raise and address the three objections most likely to be directed at
this framework.

\subsection{Objection 1: the Planck cutoff is arbitrary}

\textit{The vacuum energy calculation is sensitive to the ultraviolet
cutoff. The choice of Planck cutoff is not uniquely determined by
established physics and alternative cutoffs give completely different
results. The agreement within a factor of 1.2 is therefore cutoff-dependent
and not a robust prediction.}

\textbf{Response:} Within the IAM framework the Planck cutoff is
physically determined, not arbitrary. The calculation is asking: what
fraction of the complete Planck-scale vacuum potential has been
actualized through baryonic decoherence? The natural unit for this
fraction is the ratio of one Planck-area bit to the total horizon area
in Planck units. The Bekenstein-Hawking entropy identifies $l_P^2$ as
the area per bit of horizon information — it is the quantum of the
encoding surface. Computing the vacuum energy at any lower cutoff would
undercount the total available states and produce an incoherent ratio.
The Planck cutoff is the only physically motivated choice for this
specific calculation within IAM.

We acknowledge that this argument is specific to the IAM framework and
does not apply to standard QFT vacuum energy calculations where the
Planck cutoff is indeed a regularization choice. The prediction
equation~(\ref{eq:prediction}) is therefore a conditional result: it
holds if IAM's physical identification of the vacuum energy as the
Planck-scale potential of the horizon encoding substrate is correct.

\subsection{Objection 2: the $\Omega_b/\Omega_m$ factor was introduced
to improve agreement}

\textit{The factor $\Omega_b/\Omega_m$ was not present in the initial
formulation and appears to have been introduced post-hoc to reduce the
discrepancy from a factor of 12 to a factor of 1.2. This constitutes
parameter fitting rather than prediction.}

\textbf{Response:} The $\Omega_b/\Omega_m$ factor is not a free
parameter and was not introduced to improve numerical agreement. It is a
necessary consequence of IAM's physical identification of dark matter,
established in prior work \citep{mahaffey2026_virial_dm} independently
of the cosmological constant calculation.

In IAM, dark matter is the accumulated geometric potential half of all
prior gravitational decoherence events — it is spacetime curvature
crystallized from the potential energy of decoherence. As such, dark
matter is already classical, already actualized geometry. It does not
undergo new quantum decoherence events because it is already on the
actuality side of the quantum-to-classical transition. Only baryons —
matter that continues to exist in quantum superposition and interact
gravitationally to produce new irreversible information — write new bits
onto the horizon.

The argument proceeds in the following logical order, which reflects the
historical development:
\begin{enumerate}
\item IAM identifies dark matter as accumulated geometric potential
\citep{mahaffey2026_virial_dm}.
\item This identification implies dark matter does not independently
decohere.
\item Therefore only baryons contribute to the ongoing Landauer cost
accumulation.
\item Therefore the effective writing fraction is $\Omega_b/\Omega_m$.
\item Therefore equation~(\ref{eq:prediction}) follows.
\end{enumerate}

Step 1 precedes the cosmological constant calculation both logically and
chronologically. The factor $\Omega_b/\Omega_m$ was not chosen to fit the
data. It was derived from the framework and then applied to the
calculation.

\subsection{Objection 3: the remaining factor of 1.2 indicates an
incomplete theory}

\textit{A factor of 1.2 discrepancy indicates the calculation is not
exact. Without an explanation for this factor, the agreement may be
coincidental rather than physical.}

\textbf{Response:} The factor of 1.2 has a specific physical origin that
has now been identified and computed. The correction arises from the
departure from exact de Sitter geometry: the Gibbons--Hawking temperature
of the de Sitter encoding surface ($T_{dS} = T_H\sqrt{\Omega_\Lambda}$)
differs from the effective writing temperature ($T_H$) by the ratio
$\sqrt{\Omega_\Lambda} = 0.8274$. Applying this correction (equation
\ref{eq:prediction_corrected}) yields a prediction of
$1.141 \times 10^{-123}$ against the observed $1.14 \times 10^{-123}$
— agreement to within $0.07\%$ with zero free parameters. The factor of
1.2 is not an unexplained residual. It is the Gibbons--Hawking temperature
ratio between the encoding surface and the writing process, derivable
from $\Omega_\Lambda$ alone.

Three additional considerations support a physical rather than coincidental
interpretation.

First, the required correction exponent is $\Omega_\Lambda^{0.502}$,
indistinguishable from $\Omega_\Lambda^{1/2}$ at Planck parameter
precision. The square root arises naturally from the ratio of temperatures
(linear in $H$, which is quadratic in $\Omega_\Lambda$). This is not
tuned.

Second, the probability of accidentally recovering the observed ratio
$1.14 \times 10^{-123}$ within a factor of 2 from an expression
involving only the Planck length, the Hubble constant, and two standard
cosmological density parameters — with no free parameters — is
vanishingly small.

Third, the factor of 1.2 is orders of magnitude smaller than the
previous state of the problem. The cosmological constant problem spanned
$10^{122}$ orders of magnitude. Equation~(\ref{eq:prediction}) reduces
this to a factor of 1.2, and equation~(\ref{eq:prediction_corrected})
closes it to $0.07\%$.

% ─────────────────────────────────────────────────────────────────────────────
\section{The Cosmological Constant Is Not Constant}
\label{sec:not_constant}

A direct consequence of the IAM identification is that the cosmological
``constant'' is not strictly constant. Within IAM the accumulated
Landauer cost grows as the universe ages and more baryonic decoherence
events write information onto the horizon. The accumulation is governed
by the IAM activation function $E(a) = \exp(1-1/a)$
\citep{mahaffey2026}, which is monotonically increasing. As $E(a)$
advances, the accumulated Landauer cost increases, and therefore
$\Lambda$ increases slowly with cosmic time.

The rate of increase is set by $\mathrm{d}E/\mathrm{d}a|_{a=1} = 1$
(since $E(a) = \exp(1-1/a)$ gives $\mathrm{d}E/\mathrm{d}a = E(a)/a^2$,
and at $a=1$ this equals 1). The predicted equation of state parameter
is therefore:

\begin{equation}
w = -1 + \epsilon(a)
\label{eq:w}
\end{equation}

where $\epsilon(a) > 0$ is a small positive correction growing with
$E(a)$. The dark energy equation of state is predicted to be slightly
above $-1$ today and to have been more negative at earlier epochs when
$E(a)$ was smaller.

This prediction is qualitatively consistent with recent DESI DR2 results
\citep{desi2025}, which report a $2.8$--$4.2\sigma$ preference for
dynamical dark energy with $w > -1$. The DESI signal is interpreted
within IAM not as evidence for a new dark energy field but as the
observational consequence of the cosmological constant growing slowly as
baryonic decoherence accumulates on the horizon.

We note that this prediction — $w > -1$ from the slow growth of the
accumulated Landauer cost — is not a new addition introduced to
accommodate DESI DR2. It follows directly from the IAM activation
function, which was derived from horizon thermodynamics and established
in prior work \citep{mahaffey2026} independently of the DESI results.

% ─────────────────────────────────────────────────────────────────────────────
\section{The Full Decoherence History Integral}
\label{sec:integral}

The calculation in Section~\ref{sec:calculation} uses today's values
of $\Omega_b/\Omega_m$ and $l_H$. A more precise derivation requires
integrating the baryonic decoherence contribution over the full cosmic
history from the onset of the matter sector at electroweak symmetry
breaking ($a_\mathrm{EW} \approx 2.3 \times 10^{-15}$,
$T \approx 100$ GeV) to today ($a = 1$).

The accumulated cosmological constant at scale factor $a$ is:

\begin{equation}
\frac{\Lambda(a)}{\rho_\mathrm{vac}} \propto
\int_{a_\mathrm{EW}}^{a}
\frac{\Omega_b(a') / \Omega_\mathrm{tot}(a')}{l_H(a')^2 / l_P^2}\
\frac{\mathrm{d}a'}{a'^2 H(a')/H_0}
\label{eq:integral}
\end{equation}

where $\Omega_b(a')/\Omega_\mathrm{tot}(a')$ is the baryonic fraction
of the total energy density at scale factor $a'$, and $l_H(a') =
c/H(a')$ is the Hubble length at that epoch.

This integral correctly weights the contribution from each epoch. During
radiation domination ($a \ll a_\mathrm{eq}$), baryons are a small
fraction of total energy and the Hubble length is small. During matter
domination, baryons constitute a larger fraction of matter energy and
the Hubble length grows. During dark energy domination, the Hubble
length asymptotes to its present value.

The de Sitter temperature correction derived in
Section~\ref{sec:dscorrection} — a factor of $\sqrt{\Omega_\Lambda}$
applied to the present-day prediction — is consistent with the dominant
contribution to this integral arising from the late-time (dark-energy
dominated) epoch, where the de Sitter temperature mismatch is largest.
The full integral provides a complementary derivation of the same
correction, confirming that the $\sqrt{\Omega_\Lambda}$ result is robust.
The normalization and evaluation of equation~(\ref{eq:integral}) are the
subject of ongoing work and provide an independent cross-check of
equation~(\ref{eq:prediction_corrected}).

% ─────────────────────────────────────────────────────────────────────────────
\section{Relation to Prior Work}
\label{sec:prior}

The cosmological constant problem has been approached from multiple
directions. Weinberg \citep{weinberg1989} provided an anthropic bound on
$\Lambda$ but not a derivation of its value. Supersymmetric cancellations
reduce the discrepancy by some orders of magnitude but do not eliminate
it \citep{witten2000}. The string landscape \citep{bousso2000} provides
a framework for anthropic selection but does not predict $\Lambda$ from
first principles. Unimodular gravity \citep{weinberg1989} reformulates
the problem but does not resolve it.

The present approach differs from all of these in two respects. First, it
does not attempt to reduce the theoretical vacuum energy. The Planck-scale
vacuum energy density is accepted as the correct value of the total quantum
potential. Second, it does not invoke selection mechanisms or anthropic
arguments. The observed value of $\Lambda$ is a derived consequence of
the universe's decoherence history, computed from known physical constants
and standard cosmological parameters.

The closest prior approach in spirit is Verlinde's entropic gravity
\citep{verlinde2011} and Padmanabhan's thermodynamic interpretation of
gravity \citep{padmanabhan2010}, both of which connect gravitational
physics to horizon thermodynamics. The present work differs in that the
cosmological constant is identified specifically with the Landauer cost
of irreversible information production on timelike worldlines — a
specific physical process rather than a general thermodynamic argument —
and in that the prediction is quantitative rather than qualitative.

The IAM framework builds directly on the thermodynamic derivation of the
gravitational field equations by Jacobson \citep{jacobson1995} and its
extension to the FRW apparent horizon by Cai and Kim \citep{caikim2005}.
The additional term — the informational entropy contribution from
gravitational decoherence — is the specific modification that produces
both the IAM cosmological predictions \citep{mahaffey2026} and the
cosmological constant resolution presented here.

% ─────────────────────────────────────────────────────────────────────────────
\section{Conclusions}
\label{sec:conclusions}

We have presented a resolution of the cosmological constant problem
within the Informational Actualization Model. The resolution rests on a
single physical distinction: the QFT vacuum energy density and the
observed cosmological constant are not the same kind of quantity.

The vacuum energy density $\rho_\mathrm{vac}$ is the Planck-scale
energy density of the complete quantum substrate — the total
zero-point energy of all field configurations available as horizon
information states. The cosmological constant $\Lambda_\mathrm{obs}$
is the accumulated Landauer cost of converting Planck-scale vacuum
fluctuations to classical actuality through irreversible gravitational
decoherence on baryonic timelike worldlines over 13.8 billion years of
cosmic history.

The baseline prediction:
\begin{equation}
\frac{\Lambda_\mathrm{obs}}{\rho_\mathrm{vac}} =
\frac{2}{\pi}\left(\frac{l_P}{l_H}\right)^2 \times \frac{\Omega_b}{\Omega_m}
= 1.38 \times 10^{-123}
\end{equation}
recovers the observed $1.14 \times 10^{-123}$ within a factor of 1.2
with no free parameters. The factor of 1.2 arises from the departure
from exact de Sitter geometry: the Gibbons--Hawking temperature of the
de Sitter encoding surface ($T_{dS}$) is lower than the effective writing
temperature ($T_H$) by the ratio $\sqrt{\Omega_\Lambda}$. Applying this
correction:
\begin{equation}
\frac{\Lambda_\mathrm{obs}}{\rho_\mathrm{vac}}\bigg|_\mathrm{corrected} =
\frac{2}{\pi}\left(\frac{l_P}{l_H}\right)^2 \sqrt{\Omega_\Lambda}
\times \frac{\Omega_b}{\Omega_m}
= 1.141 \times 10^{-123}
\end{equation}
against the observed $1.14 \times 10^{-123}$ — agreement to within
$0.07\%$, zero free parameters beyond the standard cosmological model.

A secondary prediction follows: the cosmological constant is not strictly
constant but increases slowly as $E(a)$ advances, implying $w > -1$ at
late times. This is qualitatively consistent with DESI DR2
\citep{desi2025}.

A direct observational confirmation of the information writing constraint
has been obtained from a dedicated Planck 2018 MCMC chain in which the
BBN prior on $\Omega_b h^2$ was removed and $\Omega_b h^2$ was allowed
to range freely over a prior four times wider than standard. The Planck
CMB likelihood alone, with no nuclear physics input, returns
$\eta_\mathrm{implied} = (6.115 \pm 0.037) \times 10^{-10}$, consistent with
the observed baryon asymmetry $\eta_\mathrm{obs} = 6.14 \times 10^{-10}$
to within $0.36\%$. The cosmological constant and the baryon asymmetry are not
independent problems. They are two expressions of the same information
writing constraint.

The cosmological constant problem has persisted for over fifty years
because every proposed resolution has attempted to reduce the theoretical
vacuum energy to match the observed value. The IAM framework suggests
that both quantities are correct as computed. The issue is not the
magnitude of either quantity but the physical interpretation of their
ratio. The vacuum energy is the total potential of the quantum substrate.
The cosmological constant is the accumulated fraction that has been
actualized through the irreversible process of gravitational decoherence
over the age of the universe, corrected for the temperature ratio between
the de Sitter encoding surface and the effective Hubble writing process.

In a universe 13.8 billion years old, with a Hubble horizon $10^{61}$
Planck lengths across, containing baryonic matter constituting 15.6\% of
the total matter density, and with dark energy comprising 68.5\% of the
energy budget (setting the de Sitter temperature correction
$\sqrt{\Omega_\Lambda} = 0.827$), the expected fraction of actualized to
total potential is $1.141 \times 10^{-123}$. This is what is observed.

% ─────────────────────────────────────────────────────────────────────────────
\section*{Data Availability}

All MCMC chains, analysis scripts, and computational notebooks supporting
the IAM framework are publicly available at
\url{https://github.com/hmahaffeyges/IAM-Validation}
(Zenodo DOI:
\href{https://doi.org/10.5281/zenodo.18702042}{10.5281/zenodo.18702042}).

\section*{Acknowledgments}

The thermodynamic framework developed by T.\ Jacobson
\citep{jacobson1995} is the foundation upon which this work is built.
The author thanks the developers of CAMB \citep{lewis2002}, MGCAMB
\citep{wang2023}, and Cobaya \citep{torrado2021} for public code access.

% ─────────────────────────────────────────────────────────────────────────────
\bibliographystyle{plainnat}
\begin{thebibliography}{99}

\bibitem[Bekenstein(1973)]{bekenstein1973}
Bekenstein, J.\ D.\ 1973, Phys.\ Rev.\ D, 7, 2333

\bibitem[Bousso \& Polchinski(2000)]{bousso2000}
Bousso, R.\ \& Polchinski, J.\ 2000, JHEP, 06, 006

\bibitem[Cai \& Kim(2005)]{caikim2005}
Cai, R.-G.\ \& Kim, S.\ P.\ 2005, JHEP, 02, 050

\bibitem[Caldwell et al.(1998)]{caldwell1998}
Caldwell, R.\ R., Dave, R., \& Steinhardt, P.\ J.\ 1998,
Phys.\ Rev.\ Lett., 80, 1582

\bibitem[DESI Collaboration(2025)]{desi2025}
DESI Collaboration, et al.\ 2025, preprint arXiv:2503.14738

\bibitem[Gibbons \& Hawking(1977)]{gibbons1977}
Gibbons, G.\ W.\ \& Hawking, S.\ W.\ 1977, Phys.\ Rev.\ D, 15, 2738

\bibitem[Harris et al.(2020)]{harris2020}
Harris, C.\ R., et al.\ 2020, Nature, 585, 357

\bibitem[Hawking(1975)]{hawking1975}
Hawking, S.\ W.\ 1975, Commun.\ Math.\ Phys., 43, 199

\bibitem[Hunter(2007)]{hunter2007}
Hunter, J.\ D.\ 2007, Comput.\ Sci.\ Eng., 9, 90

\bibitem[Jacobson(1995)]{jacobson1995}
Jacobson, T.\ 1995, Phys.\ Rev.\ Lett., 75, 1260

\bibitem[Landauer(1961)]{landauer1961}
Landauer, R.\ 1961, IBM J.\ Res.\ Dev., 5, 183

\bibitem[Lewis \& Bridle(2002)]{lewis2002}
Lewis, A.\ \& Bridle, S.\ 2002, Phys.\ Rev.\ D, 66, 103511

\bibitem[Mahaffey(2026a)]{mahaffey2026}
Mahaffey, H.\ W.\ 2026a, preprint.
doi:\href{https://doi.org/10.5281/zenodo.18702042}{10.5281/zenodo.18702042}

\bibitem[Mahaffey(2026b)]{mahaffey2026_virial}
Mahaffey, H.\ W.\ 2026b, The Virial Partition Across Wide Range of
Physical Scales.
doi:\href{https://doi.org/10.5281/zenodo.18702042}{10.5281/zenodo.18702042}

\bibitem[Mahaffey(2026e)]{mahaffey2026_baryon}
Mahaffey, H.\ W.\ 2026e, The Baryon Asymmetry as a Derived Quantity:
CMB Evidence Without BBN Prior.
doi:\href{https://doi.org/10.5281/zenodo.18702042}{10.5281/zenodo.18702042}

\bibitem[Mahaffey(2026c)]{mahaffey2026_virial_dm}
Mahaffey, H.\ W.\ 2026c, Dark Matter and Dark Energy as Virial Partners.
doi:\href{https://doi.org/10.5281/zenodo.18702042}{10.5281/zenodo.18702042}

\bibitem[Mahaffey(2026d)]{mahaffey2026_higgs}
Mahaffey, H.\ W.\ 2026d, The Higgs Boson and the Origin of Duration.
doi:\href{https://doi.org/10.5281/zenodo.18702042}{10.5281/zenodo.18702042}

\bibitem[Martin(2012)]{martin2012}
Martin, J.\ 2012, C.\ R.\ Phys., 13, 566

\bibitem[Padmanabhan(2010)]{padmanabhan2010}
Padmanabhan, T.\ 2010, Rep.\ Prog.\ Phys., 73, 046901

\bibitem[Perlmutter et al.(1999)]{perlmutter1999}
Perlmutter, S., et al.\ 1999, ApJ, 517, 565

\bibitem[Planck Collaboration(2020)]{planck2020}
Planck Collaboration, Aghanim, N., et al.\ 2020, A\&A, 641, A6

\bibitem[Riess et al.(1998)]{riess1998}
Riess, A.\ G., et al.\ 1998, AJ, 116, 1009

\bibitem['t Hooft(1993)]{thooft1993}
't Hooft, G.\ 1993, preprint gr-qc/9310026

\bibitem[Susskind(1995)]{susskind1995}
Susskind, L.\ 1995, J.\ Math.\ Phys., 36, 6377

\bibitem[Torrado \& Lewis(2021)]{torrado2021}
Torrado, J.\ \& Lewis, A.\ 2021, JCAP, 05, 057

\bibitem[Verlinde(2011)]{verlinde2011}
Verlinde, E.\ 2011, JHEP, 04, 029

\bibitem[Wang et al.(2023)]{wang2023}
Wang, Z., et al.\ 2023, JCAP, 2023, 022

\bibitem[Weinberg(1989)]{weinberg1989}
Weinberg, S.\ 1989, Rev.\ Mod.\ Phys., 61, 1

\bibitem[Witten(2000)]{witten2000}
Witten, E.\ 2000, preprint hep-ph/0002297

\end{thebibliography}

\end{document}
